\begin{theorem}[Induction Principle for a Peano System]
\label{thm:induction-principle-for-peano-system}
Let $(P,S,1)$ be a Peano system, and let $\varphi$ be a predicate on $P$.
If
\[
\varphi(1)
\qquad\text{and}\qquad
\forall n \in P\;(\varphi(n) \Longrightarrow \varphi(S(n))),
\]
then
\[
\forall n \in P\;\varphi(n).
\]
\hyperref[prf:induction-principle-for-peano-system]{Go to proof.}
\end{theorem}
\end{theorembox}

\begin{formalizationrecord}{Lean 4 Verification Record}
\FormalizationSource{Peano, \emph{Arithmetices principia}; Feferman, \emph{The Number Systems}, Section 3.1}
\Formalizes{thm:induction-principle-for-peano-system}
\textbf{Lean module:} \texttt{LRA.VolumeII.PeanoSystems.Induction.Core}.\par
\LeanDeclaration{InductionPrincipleForPeanoSystem}
\textbf{Status:} pending against \texttt{lra-lean}.\par

\begin{leancode}
theorem InductionPrincipleForPeanoSystem
    (ps : PeanoSystem)
    (predicate : LRA.VolumeI.Set.LRASet ps.carrier)
    (base_case : predicate ps.one)
    (successor_step :
      forall element : ps.carrier,
        predicate element -> predicate (ps.successor element)) :
    forall element : ps.carrier, predicate element :=
  PeanoAxioms.PeanoInduction
    ps
    predicate
    base_case
    successor_step
\end{leancode}
\end{formalizationrecord}

\begin{remark*}[Standard quantified statement]
For a Peano system $(P,S,1)$ and a predicate $\varphi$ on $P$,
\[
\Bigl(\varphi(1) \land
\forall n \in P\;(\varphi(n) \Longrightarrow \varphi(S(n)))\Bigr)
\Longrightarrow
\forall n \in P\;\varphi(n).
\]
\end{remark*}

\begin{remark*}[Predicate reading]
\[
\operatorname{InductionPrinciple}(P,S,1).
\]
\end{remark*}

\begin{remark*}[Contrapositive quantified statement]
\[
\begin{aligned}
&\text{Let } (P,S,1) \text{ be a Peano system and let } \varphi
\text{ be a predicate on } P.\\
&\neg\forall n \in P\;\varphi(n)
\Longrightarrow
\neg\Bigl(\varphi(1) \land
\forall n \in P\;(\varphi(n) \Longrightarrow \varphi(S(n)))\Bigr).
\end{aligned}
\]
\end{remark*}

\begin{remark*}[Interpretation]
This theorem converts the subset form of induction into the predicate form
used in ordinary arithmetic arguments. A property propagates through the
entire Peano system once it holds at the distinguished first element and is
preserved by successor. The theorem therefore turns the structural induction
axiom into a proof principle for arbitrary predicates on $P$.
\end{remark*}

\begin{remark*}[Source comparison]
Peano's axiomatization includes induction as part of the natural-number
structure. Feferman derives the working form of mathematical induction from
the subset formulation of the Peano axioms immediately after Theorem 3.3 in
Chapter 3. The present notes formalize that induction principle as an explicit
reusable theorem. \citep{PeanoArithmeticesPrincipia1889}; \citet{FefermanNumberSystems1964}
\end{remark*}

\begin{dependencies}
\begin{itemize}
  \item \hyperref[def:peano-system]{Peano System}
  \item \hyperref[ax:peano-induction]{Induction Axiom}
\end{itemize}
\end{dependencies}
\end{topicbox}

\begin{topicbox}{Strong Induction}
\begin{exposition}
Ordinary induction propagates a property from each element to its immediate
successor. Strong induction allows the inductive step to use the property at
every predecessor, not just the immediate one. This broader hypothesis makes
strong induction convenient when the successor step needs information from the
entire earlier segment of the Peano system.
\end{exposition}

